% Original Prism writing draft for European Journal of Medicinal Chemistry.
% Distribution-provided elsarticle and elsarticle-num remain unchanged.
% All demonstration values and graphics are synthetic, not research findings.
\documentclass[a4paper,preprint,12pt]{elsarticle}
\usepackage{amsmath,amssymb,graphicx,array,booktabs,tabularx,siunitx,newfloat}
\usepackage[version=4]{mhchem}
\usepackage[hidelinks]{hyperref}
\providecommand{\doi}[1]{\href{https://doi.org/#1}{\nolinkurl{#1}}}
\DeclareFloatingEnvironment[fileext=los,name=Scheme,placement=htbp]{scheme}
\biboptions{numbers,sort&compress}
\setcitestyle{square}
\sisetup{uncertainty-mode=separate}
\setlength{\emergencystretch}{2em}
\raggedbottom
\clubpenalty=10000
\widowpenalty=10000
\displaywidowpenalty=10000
\renewcommand{\topfraction}{0.9}
\renewcommand{\textfraction}{0.08}
\renewcommand{\floatpagefraction}{0.8}
% Keep draft float pages compact instead of stretching space between graphics.
\makeatletter
\setlength{\@fptop}{0pt}
\setlength{\@fpsep}{18pt plus 2pt minus 2pt}
\setlength{\@fpbot}{0pt plus 1fil}
\makeatother
\journal{European Journal of Medicinal Chemistry}

\begin{document}
\begin{frontmatter}
\title{Manuscript title}
\author[inst1]{First Author}
\author[inst1,inst2]{Second Author\texorpdfstring{\corref{corresponding}}{}}
\cortext[corresponding]{Corresponding author.}
\ead{second.author@example.org}
\address[inst1]{Department, First Institution, City, Country}
\address[inst2]{Department, Second Institution, City, Country}

\begin{abstract}
Replace this paragraph with a concise, self-contained statement of the
review's subject, scope, principal themes and supported implications.
Summarize only the review actually written. Follow the current journal
instructions; this starter sets no unverified word limit. The figures,
values and equations below are synthetic formatting examples, not claims
about published evidence or an account of a literature search.
\end{abstract}
\begin{keyword}
keyword one \sep keyword two \sep keyword three
\end{keyword}
\end{frontmatter}

\section{Introduction}\label{sec:introduction}
Define the review's scope, purpose and organizing question. Explain its
relation to the actual literature supplied by the authors and identify
boundaries of coverage. Do not claim a literature search, screening process
or coverage period that the authors have not established. The public
citations here demonstrate bibliography mechanics only; retain them only
if relevant~\cite{wilkinson2016fair,peng2011reproducible}.

\section{Conceptual framework}\label{sec:framework}
Explain the concepts required by the intended readership. Distinguish
established observations, interpretations and open questions, and replace
these thematic headings with those appropriate to the review.

\section{Strategies and evidence}\label{sec:evidence}
\subsection{Chemical strategies}
Organize actual published approaches around the review's question. Attribute
structures and findings accurately, and obtain permissions when necessary.
This public placeholder supplies no conclusion about a drug or target.

Scheme~\ref{sch:sequence} demonstrates an editable vector reaction-layout
placeholder. Use the authors' original structures and preserve compound
identities consistently in the text, schemes and tables.
\begin{scheme}[htbp]
\centering
\fbox{\parbox[c][24mm][c]{0.86\linewidth}{\centering
\ce{A ->[step 1] B ->[step 2] C}\\[7pt]
Replace A, B and C with original structures and verified conditions.}}
\caption{\textbf{Synthetic scheme layout.} Letters stand for layout slots,
not chemical substances. No reaction, yield, synthesis route or experiment
is reported. Replace the drawing, conditions and legend with the authors'
verified work, or remove this demonstration.}
\label{sch:sequence}
\end{scheme}

Figure~\ref{fig:overview} demonstrates a two-panel layout. Long legends
belong to the page layout: preserve readable panel labels and account for
the full caption height when choosing the graphic size.
\begin{figure}[htbp]
\centering
\begin{minipage}[t]{0.47\linewidth}
\centering\textbf{a}\par\smallskip
\fbox{\parbox[c][27mm][c]{0.90\linewidth}{\centering
Original structure slot\\[5pt]Replace with a vector graphic}}
\end{minipage}\hfill
\begin{minipage}[t]{0.47\linewidth}
\centering\textbf{b}\par\smallskip
\fbox{\parbox[c][27mm][c]{0.90\linewidth}{\centering
Evidence graphic slot\\[5pt]Keep final-size text readable}}
\end{minipage}
\caption{\textbf{Synthetic two-panel example with an extended legend.}
(a) An original structure or concept diagram may occupy this slot.
(b) An independently prepared evidence display may occupy the second slot.
These framed regions contain no data. Define the actual symbols, units,
uncertainty and relevant conditions in the replacement legend. Preserve
panel identities across text references and graphic labels; keep axes,
legends and chemical labels readable at the final printed size. Remove
instructions and unused panels before submission.}
\label{fig:overview}
\end{figure}

Table~\ref{tab:sar} illustrates a structure--activity relationship (SAR)
table with bold compound identifiers, chemical formulae and aligned errors.
The entries are invented formatting values, not assay observations.
\begin{table}[htbp]
\centering\small
\caption{Synthetic SAR-style values for column-alignment practice only.}
\label{tab:sar}
\begin{tabularx}{\linewidth}{@{}c >{\raggedright\arraybackslash}X
 S[table-format=2.1(2)] S[table-format=2.1(2)]@{}}
\toprule
Compound & Substituent slot & {\shortstack{Illustrative endpoint 1\\/ nM}}
 & {\shortstack{Illustrative endpoint 2\\/ nM}}\\
\midrule
\textbf{1} & \ce{CH3} & 12.4 +- 0.6 & 31.0 +- 1.4\\
\textbf{2} & \ce{C2H5} & 8.1 +- 0.4 & 27.6 +- 1.2\\
\textbf{3} & \ce{F} & 15.2 +- 0.8 & 19.4 +- 0.9\\
\textbf{4} & Replace with an actual group & \multicolumn{1}{c}{---}
 & \multicolumn{1}{c}{---}\\
\bottomrule
\end{tabularx}
\par\smallskip\footnotesize\raggedright
No assay was performed. Replace all values and define the actual endpoint,
units, error measure and replicate counts. A dash is a missing-value slot,
not a zero. Formulae such as \ce{CO2} demonstrate chemical typography only.
\end{table}

Equations~\eqref{eq:ratio} and~\eqref{eq:difference} demonstrate alignment
and cross-references. Define actual quantities before using any equation.
\begin{align}
 q_i &= \frac{x_i}{x_{\mathrm{ref}}},\label{eq:ratio}\\
 \Delta q &= q_{\mathrm{A}}-q_{\mathrm{B}}.\label{eq:difference}
\end{align}
These expressions are layout examples, not a validated medicinal chemistry
metric or evidence of a scientific result.

\section{Critical comparison}\label{sec:comparison}
Compare the actual evidence and explain meaningful agreements, differences
and limitations. Account for incompatible assay conditions and uncertainty
before comparing numerical findings. Avoid treating illustrative numbers
or a graphic slot as evidence of a literature conclusion.

\section{Challenges and Outlook}\label{sec:outlook}
Identify evidence gaps, limitations and future directions grounded in the
review's cited sources. Separate feasible proposals from established facts.
The review skeleton does not imply that a systematic search, experimental
study or validation was performed.

\section*{Data availability}
Replace with the authors' accurate availability statement, repository
identifiers where applicable, and any access restrictions.
\section*{Author contributions}
Describe actual contributions using the journal's current requirements.
\section*{Declaration of competing interest}
Replace with the authors' confirmed declaration.
\section*{Acknowledgements and funding}
List confirmed acknowledgements and funding, or remove if not applicable.
\section*{Supplementary material}
Describe the actual supporting files. The included \texttt{supplementary.tex}
compiles independently with S-numbered pages and displays. A separate
\texttt{graphical\_abstract.tex} and editable \texttt{highlights.txt} are
optional drafting aids; confirm their submission requirements with the
journal before using them.

% The two public references test mechanics only; replace as appropriate.
\bibliographystyle{elsarticle-num}
\bibliography{references}
\end{document}
